%%=============================================================================
%% Samenvatting
%%=============================================================================

% TODO: De "abstract" of samenvatting is een kernachtige (~ 1 blz. voor een
% thesis) synthese van het document.
%
% Een goede abstract biedt een kernachtig antwoord op volgende vragen:
%
% 1. Waarover gaat de bachelorproef?
% 2. Waarom heb je er over geschreven?
% 3. Hoe heb je het onderzoek uitgevoerd?
% 4. Wat waren de resultaten? Wat blijkt uit je onderzoek?
% 5. Wat betekenen je resultaten? Wat is de relevantie voor het werkveld?
%
% Daarom bestaat een abstract uit volgende componenten:
%
% - inleiding + kaderen thema
% - probleemstelling
% - (centrale) onderzoeksvraag
% - onderzoeksdoelstelling
% - methodologie
% - resultaten (beperk tot de belangrijkste, relevant voor de onderzoeksvraag)
% - conclusies, aanbevelingen, beperkingen
%
% LET OP! Een samenvatting is GEEN voorwoord!

%%---------- Nederlandse samenvatting -----------------------------------------
%
% TODO: Als je je bachelorproef in het Engels schrijft, moet je eerst een
% Nederlandse samenvatting invoegen. Haal daarvoor onderstaande code uit
% commentaar.
% Wie zijn bachelorproef in het Nederlands schrijft, kan dit negeren, de inhoud
% wordt niet in het document ingevoegd.

\IfLanguageName{english}{%
\selectlanguage{dutch}
\chapter*{Samenvatting}
\lipsum[1-4]
\selectlanguage{english}
}{}

%%---------- Samenvatting -----------------------------------------------------
% De samenvatting in de hoofdtaal van het document

\chapter*{\IfLanguageName{dutch}{Samenvatting}{Abstract}}
\vspace{1em}Stress komt in elke job voor. In deze bachelorproef wordt er gekeken naar de IT-sector. Waar gemiddeld 29\% van de Europese werknemers lijdt aan stress of depressie, loopt dat cijfer bij IT-professionals op tot bijna het dubbele. Een belangrijke en onderbelichte oorzaak hiervan is een onevenwichtige taakverdeling: taken worden te vaak toegewezen op basis van intuïtie of gewoonte, zonder rekening te houden met de werkdruk of capaciteiten van werknemers.

In deze proef wordt er onderzocht of artificiële intelligentie een oplossing kan bieden voor dit probleem. De centrale onderzoeksvraag luidt dan ook: ``Hoe kan artificiële intelligentie worden ingezet om de taakverdeling binnen de IT-sector te optimaliseren, om zo enerzijds de mentale gezondheid bij werknemers te verbeteren en anderzijds de werking van het bedrijf vlotter te laten verlopen?''. Het doel is een goed onderbouwd AI-conceptmodel te ontwikkelen dat aantoont hoe een persoonlijk afgestemd takenpakket de werkdruk kan verminderen en daardoor de mentale gezondheid kan ondersteunen.

Het onderzoek werd uitgevoerd via Python met scikit-learn, waarbij synthetisch gegenereerde datasets werden gebruikt: 22 werknemers en 100 taken, elk voorzien van verschillende eigenschappen. Op basis van een heuristisch puntensysteem werd een Random Forest classifier getraind die voor elke werknemer-taakcombinatie een geschiktheidsscore voorspelt. Na de training werd het model ingezet in een simulatie waarbij de werkdruk van elke werknemer na elke toewijzing wordt bijgewerkt.

Het model werd vergeleken met Round-Robin, Greedy Rol-Match en ChatGPT (GPT-4o) op zes kwantitatieve metrics. Greedy en ChatGPT behaalden een rol-match van 100\%, maar creëerden ernstige onevenwichten in de werklastverdeling met standaarddeviaties van 3,13 en 3,35. Het ML-model behaalde een rol-match van 46\%, dat bewust lager ligt door de afweging met werkdruk en persona-fit. Het heeft wel een werklastverdeling met standaarddeviatie van slechts 1,37, waarbij elke werknemer minimaal drie en maximaal zeven taken ontving.

Het ML-model is de enige methode die inhoudelijke kwaliteit en werklastbalans systematisch combineert, wat het duurzamer maakt voor het welzijn van het team. De voornaamste beperking is de afhankelijkheid van synthetische trainingsdata, aangezien training op echte historische data noodzakelijk is. Verder onderzoek kan zich richten op validatie in een reële bedrijfsomgeving, 
de integratie van reinforcement learning en een longitudinale meting van de impact van dit systeem op stressniveaus bij IT-professionals.